% Thesis introduction. Longer than paper/sections/introduction.tex, which is not input here. No result numbers except in the summary of answers,
% which are copied from the generated tables.
\section{Motivation}
This thesis develops an explainable and hallucination-aware framework in which chest CT supervises models that read chest radiographs. It
has two parts. The first, which takes most of the thesis, trains bone suppressors on projections rendered from CT and asks whether they
erase lesions; it also asks whether classifiers gain from the suppressed images. The second uses CT to teach a Swin transformer classifier
of radiographs directly, through a CT classifier and projections of the same patients, and asks whether that classifier becomes more
accurate, whether its explanations agree with those of the CT model, and whether it reports fewer findings that are not there. Both parts
ask the same question of a CT-supervised model: what it keeps, what it invents and where it looks.

A chest radiograph projects the whole thorax onto one plane. Ribs, clavicles and the spine lie over most of the lung fields, and a small nodule
or an early consolidation that sits behind a rib has to be seen through the bone shadow. Software that suppresses rib and clavicle shadows has
been shown to help readers find lung nodules \citep{freedman2011lung}, and computer-aided detection has been evaluated on such images as well
\citep{schalekamp2014cad}.

Dual-energy subtraction (DES) separates bone and soft tissue physically: two exposures at different tube voltages are combined so that either
material cancels \citep{kuhlman2006des}. It needs dedicated hardware and a second exposure, and it is rare outside larger centres. Learning-based
bone suppression produces a comparable soft-tissue image from a single standard radiograph. Two families exist today. One trains generative
models on real dual-energy pairs \citep{chen2024bsdiff,sun2026bsldm,sun2025gllcm,sun2026debonedit}; it needs paired data from the target
hospital. The other renders paired supervision from chest CT: each scan is split into bone and soft tissue, the parts are projected separately,
and a network learns to predict one component from their sum \citep{angaitkar2026anatomy}. The second family can be trained on public data,
and the bone image of every training sample is known exactly.

CT can also supervise a radiograph classifier more directly. Public data sets now pair chest CT volumes with report-derived findings
\citep{hamamci2026ctrate}, and the projections of those volumes \citep{angaitkar2026anatomy} pair each scan with a radiograph-like image of
the same patient. A classifier trained on CT sees each finding without the overlap of ribs, heart and mediastinum, and knowledge
distillation \citep{hinton2015distilling} offers a way to pass what it learned to an X-ray model through the paired images. Swin
transformers \citep{liu2022swinv2} can serve both modalities with one architecture, which makes the transfer and the comparison of their
heat maps straightforward.

\section{Problem statement}
A suppressor is usually judged by how closely its output matches a reference image: PSNR and SSIM against a dual-energy soft-tissue image, or
the error of the bone image on rendered data. These measures reward a correct bone estimate on average over the image. They do not ask what
happens to a lesion. A network that predicts the bone image is free to assign part of a nodule's signal to that image, because the nodule is
small, dense and, behind a rib, hard to tell from bone. The suppressed image $s=x-f(x)$ then shows less of the lesion than the radiograph did,
which defeats the purpose of suppression. Fidelity metrics hardly change, because the lesion covers a few hundred pixels of an image of a
million.

Measuring this needs a lesion whose signal is known, and real radiographs do not come with one. A second difficulty is that the safest model
by this measure is one that does nothing: the identity removes no bone and no lesion. Any measure of lesion loss must therefore be read next to
a measure of how much bone the model removes. Figure~\ref{fig:audit_maps} (Chapter~\ref{ch:audit}) shows the problem on one radiograph: the
same added nodule is moved into the bone image to very different degrees by three suppressors.

Classifiers raise the mirror-image problem. A classifier can report a finding the image does not contain, or miss one it does, and its
accuracy on a test set does not say which images or which changes provoke such errors. Heat maps are meant to show where a classifier looks,
but they are often not checked against expert annotations, and the common gradient-based maps were designed for convolutional networks
rather than transformers \citep{chattopadhay2018gradcampp}. When CT is used to teach an X-ray classifier, three things therefore need testing besides accuracy: whether the transfer helps
on real radiographs at all, whether the two models attend to the same place in the same patient, and whether the classifier hallucinates
less when the image changes in ways that add anatomy but no disease.

\section{Research questions}
The thesis answers six research questions; each is tied to pre-specified hypotheses of the four analysis plans (Chapters~\ref{ch:v1},
\ref{ch:v2}, \ref{ch:classifier} and~\ref{ch:crossmodal}).
\begin{description}
\item[RQ1] Is supervision with exact CT components enough to keep lesions, and how much of a lesion does a model trained on rendered LIDC-IDRI
 data absorb? (first plan, H1; Chapter~\ref{ch:v1})
\item[RQ2] Does a lesion-invariance (twin) regulariser lower absorption without losing fidelity on real radiographs? (second plan, H1$'$, H2$'$
 and H3$'$; Chapters~\ref{ch:v2}, \ref{ch:audit} and~\ref{ch:benchmark})
\item[RQ3] Does the counterfactual audit measure something real, that is, does evidence that does not use lesion insertion agree with it?
 (H4$'$a; Chapter~\ref{ch:benchmark})
\item[RQ4] Does it matter for downstream tasks, and how reliable is the output: does the ensemble spread track the error, and do conformal
 intervals cover it, also on a population no model was trained on? (H4$'$b and H5$'$; Chapters~\ref{ch:benchmark}, \ref{ch:reliability}
 and~\ref{ch:external})
\item[RQ5] Do classifiers gain from the suppressor's soft-tissue and bone images when they read them as separate streams, compared with a
 control of the same capacity, and for which findings? (third plan, H6--H8 and H10; Chapter~\ref{ch:classifier})
\item[RQ6] Does knowledge distilled from a Swin CT classifier into a Swin X-ray classifier improve the classification of real radiographs,
 do the CT and X-ray models attend to the same place in the same patient, and does the transfer make the classifier hallucinate less?
 (fourth plan, H11--H15; Chapter~\ref{ch:crossmodal})
\end{description}

\section{Contributions}
\begin{enumerate}
\item A counterfactual lesion audit that measures the fraction of an added lesion a suppressor moves into its bone image, on real radiographs,
 at sites shared by all models, with bone overlap taken from independent annotations, and reported beside suppression strength (first paper and
 thesis).
\item A twin lesion-invariance regulariser for CT-supervised suppressors, with a dose-response over its weight and an ablation of its parts
 (first paper and thesis).
\item A training recipe on public data only (rendered LIDC-IDRI CTs in the first study, the Qure.ai CT projections in the second), with a
 three-model ensemble, split-conformal intervals and a screen of modern encoders (first paper and thesis; the v1 study, the renderer and the
 physics are described in full only here).
\item Two pre-specified evaluations on locked sets, with downstream classification and detection, an external validation on paediatric
 radiographs, and every change to the analysis disclosed, including negative results (first paper and thesis; the negative results and the
 deviations log are thesis appendices).
\item A pre-specified test of whether classifiers gain from the soft-tissue and bone images, used alone, stacked or as streams of a multi-stream
 network, against a control of the same capacity, with an analysis by finding group, stream-wise attribution maps scored against expert boxes
 and conformal uncertainty (second paper and thesis).
\item An explainable, hallucination-aware Swin transformer framework for cross-modality chest CT and X-ray classification: a 2.5D Swin CT
 classifier on CT-RATE, its distillation into an X-ray classifier through paired projections, a test of heat-map agreement between the
 modalities, counterfactual hallucination tests and conformal false-negative control, evaluated under a fourth pre-specified plan against a
 DenseNet-121 reference (thesis).
\end{enumerate}

\section{Outline of the thesis}
Chapter~\ref{ch:background} introduces chest radiography, dual-energy subtraction, learning-based suppression, the architectures, metrics,
uncertainty methods and statistics used later. Chapter~\ref{ch:data} describes the data sets and the CT-to-radiograph physics of the first
study. Chapter~\ref{ch:v1} reports the first, pre-specified study on rendered LIDC-IDRI data and the real-image check that led to the second
plan. Chapter~\ref{ch:v2} presents the Qure.ai projections as an exactly additive training source, the twin regulariser, the exploratory
screening that fixed its operating point, the dose-response and the encoder screen. Chapter~\ref{ch:audit} defines and validates the audit and
reports the confirmatory audit. Chapter~\ref{ch:benchmark} gives the experimental design and the locked-set results for fidelity, rib
suppression, real nodules and downstream tasks. Chapter~\ref{ch:reliability} reports the ensemble uncertainty, and Chapter~\ref{ch:external}
the external validation on BDCXR-3257. Chapter~\ref{ch:classifier} tests whether classifiers gain from the suppressor's images (third plan),
and Chapter~\ref{ch:crossmodal} builds and tests the cross-modality Swin framework (fourth plan). Chapter~\ref{ch:discussion} discusses the
findings of the suppression studies and their limits, Chapter~\ref{ch:ethics} covers reproducibility, licences and ethics, and
Chapter~\ref{ch:conclusion} answers the research questions. Appendix~\ref{app:deviations} lists the deviations from the analysis plans,
Appendix~\ref{app:negative} the negative and null results, Appendix~\ref{app:extra} additional tables, and Appendix~\ref{app:classifier}
the deviations and additional figures of the classifier study. Two papers accompany this thesis: the first reports the second study, the
second the classifier study of Chapter~\ref{ch:classifier}. The v1 study, the physics details, the screening history, the cross-modality
study, the negative results and the ethics chapter are given in full only here.

\section{Notation}
\begin{center}
\begin{tabular}{@{}lp{0.68\linewidth}@{}}
\toprule
Symbol & Meaning \\
\midrule
$x$ & radiograph (or training composite), intensities in $[0,1]$ \\
$b$, $s=x-b$ & bone image and soft-tissue image \\
$f$, $f_k$ & suppressor (predicts the bone image), member $k$ of an ensemble of $K$ \\
$\bar f$, $\sigma$ & ensemble mean and per-pixel standard deviation (spread) \\
$n$, $\alpha$ & lesion signal added to an image, its amplitude \\
$A$ & absorption: fraction of $\alpha n$ moved into the bone estimate (Equation~\ref{eq:absorption}) \\
$w$ & weight of the twin lesion-invariance term (Equation~\ref{eq:twin}) \\
$M$ & lesion support used by the twin term \\
$X$, $S$, $\hat B$ & classifier inputs in Chapter~\ref{ch:classifier}: radiograph, predicted soft-tissue and bone image \\
P0, P1 & projection models of Chapter~\ref{ch:crossmodal}, without and with distillation from the CT teacher \\
F$_\mathrm{in}$, F$_\mathrm{p0}$, F$_\mathrm{p1}$ & NIH classifiers of Chapter~\ref{ch:crossmodal}, initialised from ImageNet, P0 and P1 \\
\bottomrule
\end{tabular}
\end{center}
